% !PRINTED 297
\weberpage{297}
\begin{center}
\S.\,75.\\
\gesperrt{Metacyclic equations whose degree is a power of a prime.}
\end{center}

After the last theorems, the interest of further investigations on metacyclic equations concentrates chiefly on the case in which the degree of the equation is a power $p^k$ of a prime $p$. We assume the equation to be irreducible, and restrict ourselves to the consideration of \gesperrt{primitive equations}; for the solution of the imprimitive ones reduces to the successive solution of two (or more) primitive equations, whose degrees are likewise powers of $p$, and which, if the original equation is metacyclic, must also be metacyclic.

We now assume, then, that $P$ is the group of an irreducible primitive metacyclic equation $f(x) = 0$ of degree $p^k$. By Bd. I, \S.\,158, 2., not only $P$ itself, but all its normal divisors (with the exception of the unit group) must still be \gesperrt{transitive}. We now apply Theorem IV., proved in general in \S.\,8, according to which $P$ possesses a normal divisor $Q$ consisting entirely of commutative elements; and if several such divisors are present, we understand by $Q$ one of them whose degree is as low as possible, but still greater than 1. This group $Q$ is therefore likewise still transitive. But $Q$ can have no further proper divisor different from unity that is at the same time a normal divisor of $P$, because otherwise this divisor would take the place of $Q$.

When, by the appropriate adjunctions, the group of our equation has been reduced from $P$ to $Q$, then $f(x) = 0$ has become an \gesperrt{Abelian equation} in the extended domain of rationality, and since it has still remained irreducible, by Bd. I, \S.\,162 the degree of the equation is equal to the degree of the group; i.\,e. $Q$ is an \gesperrt{Abelian group} of degree $p^k$. The degrees of all the elements of $Q$ are therefore powers of $p$. We wish to prove that, apart from the unit element, \gesperrt{only elements of degree $p$ itself occur in $Q$.}

Let us assume that $p^\lambda$ is the highest degree that occurs among the elements of $Q$, and that $\lambda > 1$. Then all the elements of $Q$ whose degree is a divisor of $p^{\lambda-1}$ form a
%CONT>>
